\section{Conclusion}
\label{sec:conclusion}

We presented GaitNet, a footstep planner that is given no gait and
that separates choosing a step from executing it. GaitNet picks one
step per tick from sampled safe footholds, and an execution layer
makes each controller land that step where and when GaitNet meant.
Trained on one controller and run with its weights frozen on OCS2,
GaitNet succeeds about as often as OCS2's own trot on held-out worlds
with exact terrain, at the lower of our two speeds. In our pair of
controllers, the obstacles we found lay in execution rather than in
the choice of step: correcting only the landing point recovers an
earlier planner at speed on its own controller, and on OCS2 the
largest obstacle was when steps were delivered. Separating what to
step on from how a controller lands the foot there thus appears to be
a practical way to reuse a learned footstep planner across
controllers.

\section{Limitations}
\label{sec:limitations}

\pending{All results are in simulation; GaitNet has not yet run on a
Unitree Go1.} Within simulation, we moved GaitNet between only one pair
of controllers, and the OCS2 comparison uses \corlGzTrialsPerArm{}
trials per condition, which cannot separate GaitNet from OCS2's own
trot. The execution layer was also tuned with the same policy, and GaitNet
did worse than OCS2's own trot on flat ground at the faster speed
(Sec.~\ref{subsec:exp-transfer}). Nor have we ablated the execution layer on OCS2: we have
not run GaitNet there without anchoring, or with a simple hand-made
correction of the landing shift in place of the learned IDM, which
would show how much of the gain needs learning at all. At
\corlGzSpeed\,m/s the landing shift is also small, so our OCS2 tests
do not stress it. \pending{The OCS2 results reported here still use the training
controller's IDM.} Refitting the IDM requires logging steps on each new
controller first; adapting it online, in the spirit of
\cite{xie_fada_2026}, could remove that step.

The evidence for GaitNet's design is also incomplete. The value of the
candidate mixer and the history is not isolated: with the IDM, the
per-candidate network is within noise of GaitNet in Isaac Lab, it is
more robust on an OCS2-like executor on flat ground, and it has not
been tested on OCS2 from a healthy checkpoint. The deployed
checkpoint comes from an unseeded run whose settings changed during
training, and seeded runs of the final setting are still training.

Finally, our tests are mild in several ways. The success rule, half the
commanded distance (Sec.~\ref{subsec:exp-setup}), is lenient.
\pending{We will add the distance walked and a stricter threshold.}
Commands stay at or below \corlCmdVxMax\,m/s, and the weakest settings
are 50\% voids at the faster speed in Isaac Lab and the faster speed
with the camera on OCS2. GaitNet may pick any leg pair but pairs its
legs diagonally, and this freedom did not raise success here; harder
terrain or faster commands may need other pairs, which we have not
shown. GaitNet also chooses from a rule-based mask of safe cells, and
reading depth or height maps directly could help where the rules are
too strict or too loose.
